\chapter{Time Evolution in Quantum Mechanics}\label{sec:time_evolution}
Quantum dynamics describes the evolution of quantum states $\ket{\Psi(t)}\in\mathcal{H}$ over time and is fundamental in quantum chemistry for understanding molecular behavior at the atomic level. It is governed by the principles of quantum mechanics, particularly the \acrfull{tdse}
\begin{equation}\label{eq:tdse}
i\frac{\symrm{d}}{\symrm{d}t}\ket{\Psi\left(t\right)}=H\ket{\Psi\left(t\right)},
\end{equation}
which dictates how wavefunctions change with time. Due to the complex nature of these equations, practical solutions are often limited to small or model systems.

This section describes, how we can practically obtain the solution to the \acrshort{tdse}. We work in the complex separable Hilbert space $\mathcal{H}$ introduced in Sec~\ref{sec:hilbert_spaces_and_quantum_states}, with self-adjoint Hamiltonian $H$ (Sec~\ref{sec:self_adjoint_operators_and_observables}), so that the time-evolution operators described in this chapter are consistent with functional calculus treatment in Sec~\ref{sec:spectral_theory_in_hilbert_spaces}.
\input{chapters/02-time_evolution_in_quantum_mechanics/sections/01-numerical_schemes_for_tdse}
\section{Holy Trinity of Variational Principles}\label{sec:variational_principles}
\subsection{Dirac--Frenkel Principle}\label{sec:dirac_frenkel_principle}
\subsection{McLachlan's Variational Principle}\label{sec:mclachlan_variational_principle}
\subsection{Time-Dependent Variational Principle}\label{sec:time_dependent_variational_principle}

